% Einheit 2 von 5 – Dezimalzahlen addieren und subtrahieren (bank/bruchrechnung/e2.jsonl, zusammenbau v0.1)
%% TODO Zweigzeile Teil 1: was hier gelernt wird (Satz in Schülersprache aus dem Einheitstitel) – Einheit 2
\einheitenkopf[e2]{Einheit 2 von 5 · Dezimalzahlen addieren und subtrahieren}
\zweigzeile{ab Kl. 5, je nach Buch bis Kl. 6 (am Gymnasium ab Kl. 6) · keine P10-Aufgabe}
\setcounter{aufgabe}{13}

%% TODO Titel: Ich-kann-Satz fehlt in der Bank (Platzhalter: Kettenname) – Nr. 14
%% TODO Anweisung: ein Satz über der ersten Teilaufgabe fehlt in der Bank – Nr. 14
\begin{aufgabe}{Dezimal addieren/subtrahieren}
%% TODO \rechenplatz{n}: Schrittzahl des Grundfalls fehlt in der Bank (nur bei fester Schreibform, 2.3 b)
\begin{teile}
\teil Welche Zahl steht in der Stellenwerttafel?

\sachtabelle{ccc}{Einer & Zehntel & Hundertstel}{4 & 3 & 6}
\leerfeld
\teil $3{,}4 + 2{,}5$
\teil $5{,}6 + 1{,}23$
\teil $2{,}68 + 3{,}57$
\teil $3{,}45\text{ €} + 1{,}80\text{ €}$
\teil Familie Aydin hat im letzten Jahr $1\,847{,}60$ € für Heizöl bezahlt, in diesem Jahr $2\,016{,}35$ €. Zeige, dass sie in diesem Jahr $168{,}75$ € mehr bezahlt hat. \hfill (P10 2014 OS)
\end{teile}
\end{aufgabe}

%% TODO Titel: Ich-kann-Satz fehlt in der Bank (Platzhalter: Kettenname) – Nr. 15
%% TODO Anweisung: ein Satz über der ersten Teilaufgabe fehlt in der Bank – Nr. 15
\begin{aufgabe}{ganze Zahl plus Dezimalzahl}
\begin{teile}
\teil $4 + 2{,}35$
\end{teile}
\end{aufgabe}

%% TODO Titel: Ich-kann-Satz fehlt in der Bank (Platzhalter: Kettenname) – Nr. 16
%% TODO Anweisung: ein Satz über der ersten Teilaufgabe fehlt in der Bank – Nr. 16
\begin{aufgabe}{Überschlag vorab}
\begin{teile}
\teil Erst überschlagen, dann rechnen: $4{,}9 + 3{,}2$
\end{teile}
\end{aufgabe}

%% TODO Titel: Ich-kann-Satz fehlt in der Bank (Platzhalter: Kettenname) – Nr. 17
%% TODO Anweisung: ein Satz über der ersten Teilaufgabe fehlt in der Bank – Nr. 17
\begin{aufgabe}{Dezimal addieren/subtrahieren – Fehler finden, Begründen, Anwendung}
\begin{teile}
\teil Paul rechnet: $3{,}5 + 0{,}42 = 0{,}77$. Finde den Fehler.
\teil Warum schreibt man beim Addieren Komma unter Komma?
\teil Auf dem Kassenzettel stehen Brot 2,79 €, Käse 3,45 € und Milch 1,19 €. Du zahlst mit einem 10-€-Schein. Wie viel Rückgeld bekommst du?
\end{teile}
\end{aufgabe}
